% ==============================================================================
% 3. REFERENCIAL TEÓRICO
% ==============================================================================
\chapter{Referencial teórico}
\label{cap:referencial}

\section{Cobertura e suas limitações}

Critérios de cobertura medem quanto do programa uma suíte exercita: linhas, ramos, condições \cite{ammann2016}. São baratos de medir e fáceis de comunicar, e por isso viraram a métrica padrão de ``qualidade de testes'' em ferramentas, em equipes e em boa parte da pesquisa. O problema é que cobertura mede execução, não verificação. Um teste que chama uma função e não afirma nada sobre o resultado cobre as mesmas linhas que um teste que afirma o valor exato. \citeonline{inozemtseva2014} mediram, em projetos Java de tamanhos variados, a correlação entre cobertura e a capacidade da suíte de detectar defeitos, e encontraram correlação fraca a moderada, que desaparece quando se controla pelo número de testes. A conclusão deles é direta: cobertura indica o que \textit{não} foi testado, mas uma cobertura alta não é evidência de que o que foi testado está bem testado.

\section{Teste de mutação}

Teste de mutação inverte a pergunta. Em vez de medir o que a suíte executa, mede o que ela detecta. A ideia, formulada por \citeonline{demillo1978}, é introduzir pequenas alterações sintáticas no programa (os \textit{mutantes}: trocar \texttt{<} por \texttt{<=}, negar uma condição, remover uma instrução) e rodar a suíte contra cada versão alterada. Se algum teste falha, o mutante foi \textit{morto}; se todos passam, ele \textit{sobreviveu}, e a suíte não distingue o programa original daquela variante. O escore de mutação é a fração de mutantes mortos. A hipótese do programador competente e o efeito de acoplamento, discutidos em detalhe na revisão de \citeonline{papadakis2019}, justificam por que alterações pequenas servem como substitutos de defeitos reais, e \citeonline{just2014} deram suporte empírico a isso ao mostrar, com 357 defeitos reais, que suítes que matam mais mutantes também detectam mais defeitos reais, de forma independente da cobertura.

Dois detalhes práticos importam para este trabalho. O primeiro são os \textbf{mutantes equivalentes}: alterações que não mudam o comportamento observável do programa (por exemplo, trocar o valor inicial de uma variável que é sempre reatribuída antes de ser lida). Nenhum teste consegue matá-los, e a decisão de equivalência é, no caso geral, indecidível \cite{papadakis2019}; ferramentas não os identificam, e eles ficam no denominador, rebaixando o escore máximo atingível. O segundo são os \textbf{\textit{timeouts}}: quando um mutante faz o programa não terminar (um laço cuja condição de parada foi alterada), a ferramenta interrompe a execução após um limite de tempo e conta o mutante como detectado. Esse limite é calculado a partir do tempo de uma execução de referência, e por isso é sensível à carga da máquina. O StrykerJS \cite{stryker}, usado aqui, implementa operadores para JavaScript e TypeScript, integra-se ao Vitest \cite{vitest} e reporta separadamente mutantes mortos, sobreviventes, por \textit{timeout} e sem cobertura.

\section{\textit{Test smells}}

Testes também são código, e código que vai ser mantido. \citeonline{vandeursen2001} cunharam o termo \textit{test smell} para padrões em código de teste que dificultam entender, manter e confiar nos testes: o \textit{Assertion Roulette} (várias asserções sem mensagem num mesmo teste, de modo que não se sabe qual falhou), o \textit{Eager Test} (um teste que verifica vários comportamentos de uma vez), o \textit{Mystery Guest} (dependência de recursos externos ao teste), entre outros. \citeonline{meszaros2007} sistematizou esses padrões e os remédios correspondentes. Para JavaScript e TypeScript, o SNUTS.js \cite{oliveira2024snuts} detecta quinze \textit{smells} por análise da árvore sintática do arquivo de teste, o que permite medir manutenibilidade de forma automática e reprodutível. \citeonline{siddiq2024} encontraram \textit{smells} como asserções duplicadas e testes vazios em suítes geradas por LLMs, o que motiva incluir essa dimensão na avaliação.

\section{Geração de testes com LLMs e suas linhas de base}

A geração automática de testes antecede as LLMs: ferramentas baseadas em busca, como o EvoSuite, maximizam cobertura por algoritmos evolutivos, e \citeonline{lemieux2023} mostraram que uma LLM pode ajudar essas ferramentas a escapar de platôs de cobertura. Com modelos de propósito geral, a geração direta a partir de um prompt se tornou viável. \citeonline{schafer2024} avaliaram essa abordagem em 25 pacotes JavaScript e obtiveram cobertura mediana de instruções de 70\%, com uma fração relevante de testes que falhavam ou não compilavam. \citeonline{yuan2024} analisaram testes gerados pelo ChatGPT para Java e encontraram cobertura comparável à de ferramentas tradicionais, mas problemas frequentes de correção, que motivaram um processo iterativo de correção guiado por feedback do compilador e do executor. \citeonline{siddiq2024} chegaram a conclusões semelhantes para três modelos de código.

Há duas lacunas nesses estudos que este trabalho tenta fechar. A primeira é a \textbf{linha de base}. Quando há comparação, ela é feita contra testes escritos por humanos, tomados como padrão-ouro: o \textit{benchmark} TestGenEval \cite{jain2025testgeneval}, por exemplo, mede cobertura e escore de mutação dos testes gerados em relação aos testes humanos do próprio repositório, e \citeonline{walczak2025} usam a suíte de um praticante com 44\% de escore de mutação. Suítes humanas são úteis como referência do que se faz na prática, mas não do que é possível fazer; concluir que uma LLM ``supera o humano'' contra essa referência não diz quanto falta para uma suíte realmente boa. A segunda lacuna é o \textbf{oráculo}: \citeonline{konstantinou2024} mostram que oráculos gerados por LLMs tendem a capturar o comportamento real do código, e não o pretendido, o que limita a detecção de defeitos lógicos. O escore de mutação é a métrica certa para expor isso, porque mutantes são justamente desvios do comportamento real que uma asserção frouxa deixa passar. \citeonline{wang2025mutgen} usam mutação dentro do laço de geração, realimentando a LLM com os mutantes vivos; o uso aqui é diferente: a mutação é o juiz externo, não parte da geração, e a referência é um teto maximizado, não uma suíte humana.

\section{Não determinismo}

LLMs amostram tokens de uma distribuição de probabilidade; a mesma entrada gera saídas diferentes a cada chamada, e isso não é um detalhe de configuração. \citeonline{ouyang2025} estudaram o fenômeno em geração de código com 829 problemas e observaram que, na configuração padrão, a fração de tarefas em que nenhuma das execuções repetidas produz a mesma saída de teste foi de 75,8\%, 51,0\% e 47,6\% em três \textit{benchmarks}, e que fixar a temperatura em zero reduz, mas não elimina, a variação. Os autores alertam que isso é uma ameaça à validade de conclusão de qualquer estudo que compare modelos a partir de uma única execução. Para avaliação de testes gerados, a implicação é concreta: um escore de mutação medido uma vez é uma amostra de tamanho um de uma variável aleatória cuja dispersão ninguém mediu. A QP4 deste trabalho mede essa dispersão.
